% from main.tex
%%%%%%%%%%  scan idx 388  |  folio 375  %%%%%%%%%%
\textbf{Lemme 1.10.} Avec les notations et les hypothèses de 1.1, soient \(F\) et \(G\) deux faisceaux localement libres de type fini sur \(X\), tels que :
\[
  \cl_{K^{\bullet}(X)}(F)=\cl_{K^{\bullet}(X)}(G);
\]
alors il existe un entier \(n_0\), tel que pour \(n\geq n_0\), on ait :
\[
\begin{array}{ll}
\text{i)}& f_*(F(n)),\ f_*(G(n))\text{ sont localement libres de type fini sur }S;\\[1mm]
\text{ii)}& \cl_{K^{\bullet}(S)}\bigl(f_*(F(n))\bigr)=\cl_{K^{\bullet}(S)}\bigl(f_*(G(n))\bigr).
\end{array}
\]

Soit \(A\) le groupe abélien libre engendré par les classes (à isomorphisme près) de faisceaux localement libres de type fini sur \(X\). Dire que, dans \(K^{\bullet}(X)\), \(F\) et \(G\) ont même classe, c'est dire, d'après la construction du \(K^{\bullet}\), qu'il existe des suites exactes de \(\mathcal O_X\)-Modules localement libres de type fini :
\[
  0\longrightarrow L_i\longrightarrow M_i\longrightarrow N_i\longrightarrow 0,
  \qquad i=1,\ldots,k,
\]
et des entiers \(n_i\), tels que dans \(A\) on ait l'égalité :
\[
  (F)-(G)=\sum_i n_i\bigl((L_i)-(M_i)+(N_i)\bigr).
\]
Pour tout \(n\), la tensorisation par \(\mathcal O_X(n)\) définit un automorphisme de \(A\), d'où :
\[
  (F(n))-(G(n))=\sum_i \bigl((L_i(n))-(M_i(n))+(N_i(n))\bigr),
\]
les suites \(0\to L_i(n)\to M_i(n)\to N_i(n)\to0\) restant exactes.

Soient \(B\) le groupe abélien libre engendré par les classes (à isomorphisme près) de faisceaux quasi-cohérents sur \(S\), et \(A'\subset B\) le sous-groupe engendré par les classes de faisceaux localement libres de type fini sur \(S\). Par image directe, on définit un homomorphisme \(A\to B\), et on a donc encore :
